This subsection specifies the planned disaggregation of primary outcomes by topical domain and by metric, and it documents the analysis procedures that will be used to compare the graph-guided prompting pipeline to the linear baseline within each domain. Because the provided evidence package contains no domain-specific run artifacts or per-domain summaries, the text below focuses on analysis design, metric definitions, and visualization conventions rather than on numeric results.\footnote{General background note: the current manuscript draft did not include domain-specific run artifacts in the provided evidence pack; quantitative per-domain results are reported in sections of the paper that rely on experimental run outputs.}

Metrics and per-domain aggregation. For each topical domain we compute three primary per-run metrics: (1) citation coverage, defined as the fraction of non-trivial factual claims in generated sections that are annotated with at least one source pointer; (2) hallucination rate, defined via reference-based fact-checking as the fraction of asserted factual claims for which no supporting item can be found in the curated evidence set; and (3) reviewer-style score, the mean human-assigned quality rating for sections produced for that domain. (General background knowledge: these metric definitions follow common practice for citation-aware generation evaluation and are stated here to fix terminology for per-domain comparisons.)

Per-domain summaries report measures of central tendency and dispersion for each metric: mean ± standard deviation, median and interquartile range, and the full empirical distribution visualized with a violin or empirical cumulative distribution function (ECDF) plot. Where informative, we also report raw counts (number of generated sections, number of claims assessed) alongside proportional metrics to aid interpretation of sampling variability. When reporting aggregate citation-like statistics we note the potential sensitivity of summary choices (e.g., arithmetic mean versus geometric mean) and associated confidence-interval behaviour; prior work recommends care in choosing summary statistics appropriate to the distributional properties of citation-like data \cite{web_thelwall_2015_the_precision_of_the_arithmetic_mean_geometric_m}.

Statistical comparison within domains. To evaluate whether the graph-guided pipeline yields domain-specific improvements relative to the linear baseline, we apply paired or matched tests at the section level when the experimental design produces matched pairs (for example, the same section specification generated under both methods). For approximately normally distributed per-section differences we report paired t-tests; when distributional assumptions fail we use the Wilcoxon signed-rank test. We accompany p-values with effect-size estimates (Cohen's d for parametric comparisons, rank-biserial correlation for nonparametric tests) and 95% confidence intervals. To limit the false discovery rate across the three primary metrics we apply an appropriate multiple-comparison correction (e.g., Benjamini–Hochberg). (General background knowledge: these choices reflect standard statistical practice for small- to medium-sized empirical comparisons; specific test selection is reported alongside results.)

Domain × method interaction analysis. To assess whether method effects differ across domains we fit a two-way model with fixed factors for method and domain and an interaction term. When the data structure requires it (for example, repeated annotations by the same reviewer or multiple sections sharing a template) we employ mixed-effects models with random intercepts for section template and for annotator to account for non-independence. Interaction terms are interpreted conservatively: a statistically significant interaction indicates that the magnitude (or direction) of the method effect depends on domain, whereas non-significance is not taken as proof of equivalence and must be considered together with effect-size estimates and power analyses. When per-domain sample sizes are limited we report uncertainty explicitly and avoid overinterpreting small differences \cite{web_gros_2025_per_domain_generalizing_policies_on_validation_i}.

Visualization conventions and reporting. Per-domain figures include: (a) side-by-side bar charts with error bars (mean ± SE) for quick method-by-domain comparison; (b) violin plots or ECDFs showing the full distribution and tail behaviour; and (c) difference plots (graph-guided minus linear baseline) with bootstrap-derived 95% confidence intervals to highlight directional effects. Tabular reports provide numeric summaries (mean ± sd, median [IQR], sample size) and p-values adjusted for multiple comparisons; every figure and table explicitly displays sample sizes and, where applicable, effect-size metrics. Qualitative captions identify representative examples and note which sections were paired for comparison.

Qualitative per-domain analysis. For each domain we select representative section pairs (graph-guided vs. baseline) that illustrate characteristic differences in citation alignment and claim conservatism. Each example is accompanied by: the JSON-constrained section output (truncated for brevity), the list of attached evidence pointers, and a brief annotator note describing the primary discrepancy (e.g., missing citation, unsupported claim, or improved evidence alignment). These qualitative vignettes are intended to complement the quantitative breakdown and to make failure modes interpretable in the domain context.

Limitations of the per-domain breakdown. The present write-up is an analysis plan and does not include the per-domain numerical comparisons in this section. More generally, conclusions about between-domain differences require adequate sample sizes per domain and reliable retrieval pools; limited per-domain data reduces power to detect interactions and increases uncertainty in effect estimates. In addition, summaries that rely on citation-derived counts or scores should be interpreted cautiously because citation-style data can exhibit considerable heterogeneity and be sensitive to the choice of summary statistic and confidence-interval formulae \cite{web_varin_2013_statistical_modelling_of_citation_exchange_betwee,web_adler_2009_citation_statistics,web_thelwall_2015_the_precision_of_the_arithmetic_mean_geometric_m}. We therefore report both central summaries and distributional visualizations, and we make raw annotations and analysis code available to facilitate exact replication and re-analysis.

Summary. This subsection defines how per-domain results will be computed, compared, and visualized: (i) consistent metric definitions for citation coverage, hallucination rate, and reviewer score; (ii) statistical procedures for within-domain and interaction analyses, with effect sizes and multiple-comparison control; (iii) visualization and qualitative-example conventions; and (iv) explicit caveats about limited per-domain sample sizes and the need to consider effect sizes jointly with p-values when interpreting domain-specific claims.